\documentclass[a4paper,10pt]{article}
\usepackage[margin=18mm]{geometry}
\usepackage{graphicx}
\title{EPC2361: Separate Turn-On and Turn-Off Gate Resistance}
\author{LossCalculation-SPB}
\begin{document}
\maketitle
\section{Gate-drive model}
The DUT now has independent external $R_{g,on}$ and $R_{g,off}$. An ideal
command-selected impedance is used: while the driver command exceeds 2.5 V,
the external turn-on resistance plus driver source resistance is active;
otherwise the external turn-off resistance plus driver sink resistance is active.
The command still has its finite 0--5 V edge. Internal model gate resistance
of 0.4 ohm and gate-loop inductance are additional.
This represents an abstract source/sink driver with separate output paths.
It is not a steering-diode implementation, driver-IC model, Miller clamp,
or current-limited gate supply. A physical circuit must be matched to this
assumption, including diode drops and parasitics if steering diodes are used.
The complementary device's gate remains off through its original resistance.
\section{Sweep and interpretation}
\begin{center}\begin{tabular}{ll}
@SETUP@
\end{tabular}\end{center}
The sweep contains @COUNT@ single-device double-pulse cases. First turn-off
and second turn-on use separately measured actual load currents. Target
currents are not exact because finite pulse delays affect the first ramp.
The existing energy definitions are retained: excess DUT channel/access-resistor
heat, with DC ordinary conduction subtracted for turn-on. Complementary reverse
conduction and gate-drive losses are excluded from this energy pair.
No updated full-module efficiency claim is made from this table.
\section{Pairs passing every sampled target current}
@PAIRCOUNT@ voltage/resistance/inductance combinations pass all target-current
screens. Passing requires both devices below the 90 V design ceiling and
package gate voltages within -4 to +6 V over the characterized interval.
The 100 V absolute drain rating is also recorded. Resistance values are
sampled candidates, not continuous intervals or guaranteed hardware limits.
\begin{center}\small\begin{tabular}{rrrrrr}
Bus V & Lbus nH & Rg,on ohm & Rg,off ohm & Worst Vds V & Worst energy uJ \\
\hline
@PASSING@
\end{tabular}\end{center}
A larger resistor does not automatically pass every limit. Ringing depends
on both drive paths and parasitics, while slower switching changes loss and
timing. The table's energy ranking is for DUT switching heat only.
\section{Numerical checks}
One candidate is selected for a 0.2/0.1/0.05 ns timestep check and
100/150/200 ns integration-window check. An equal-resistance split path is
also compared with the original fixed resistor at identical conditions.
@CHECKS@
These checks establish numerical consistency at selected points, not
physical accuracy or qualification of the full grid. The vendor model has
no destructive breakdown behavior. High-voltage failed cases remain in the
CSV and figures. Parallel-bank gate impedance and current sharing are not
represented by this single-device bench.
@PLOTS@
\end{document}
